\chapter{Atividade 3 -- Processos, threads, concorrência e compartilhamento de memória}
\label{chap:atvd03}

\section{Enunciado da tarefa}
\label{sec:atvd03-enunciado}

\begin{quote}
Analise e execute os códigos disponibilizados e explique o comportamento observado nos valores de saída da variável \texttt{contador}.

Discuta o que acontece em cada código, identificando possíveis problemas, causas dos comportamentos observados e uma solução adequada para cada caso.

Em seguida, discuta os conceitos de processo e thread, destacando as principais diferenças entre eles, especialmente em relação a compartilhamento de memória e execução concorrente.
\end{quote}

\section{Objetivos e requisitos}
\label{sec:atvd03-objetivos}

A atividade teve como objetivo comparar, por meio de dois programas curtos em C, o comportamento de uma variável global quando a execução concorrente é realizada com processos e quando é realizada com threads.

O elemento observado nos dois casos foi a variável global \texttt{contador}. Em ambos os programas, são criadas quatro unidades de execução, e cada uma realiza um milhão de incrementos.

Os objetivos específicos considerados foram:

\begin{itemize}
    \item analisar o comportamento de \texttt{contador} quando processos filhos são criados com \texttt{fork()};
    \item analisar o comportamento de \texttt{contador} quando threads são criadas com \texttt{pthread\_create()};
    \item explicar por que os incrementos realizados pelos processos filhos não alteram o contador do processo pai;
    \item explicar por que as threads acessam a mesma variável global;
    \item identificar a condição de corrida existente no programa com threads;
    \item distinguir isolamento de memória, compartilhamento de memória e execução concorrente;
    \item propor soluções adequadas para cada caso, sem alterar a finalidade didática dos códigos fornecidos.
\end{itemize}

Os códigos analisados foram mantidos como programas demonstrativos. Assim, o objetivo principal não foi otimizar desempenho, mas compreender o modelo de memória associado a processos e threads.

\section{Fundamentação conceitual}
\label{sec:atvd03-fundamentacao}

\subsection{Processos}

Um processo é uma instância de um programa em execução. Cada processo possui seu próprio espaço de endereçamento virtual, contendo, entre outras regiões:

\begin{itemize}
    \item segmento de código;
    \item variáveis globais;
    \item heap;
    \item pilha de execução;
    \item descritores de arquivos e demais recursos associados pelo sistema operacional.
\end{itemize}

Quando um processo é criado por meio de \texttt{fork()}, o processo filho inicia como uma cópia lógica do processo pai.

Essa cópia inclui o valor das variáveis no instante da criação do filho. Entretanto, pai e filho não passam a compartilhar automaticamente as mesmas posições de memória. Após o \texttt{fork()}, cada processo executa em seu próprio espaço de endereçamento.

Consequentemente, se um processo filho modifica uma variável global, essa modificação ocorre em sua própria cópia da variável. O processo pai não vê essa alteração.

\subsection{Threads}

Uma thread é uma linha de execução dentro de um processo.

Diferentemente dos processos, threads pertencentes ao mesmo processo compartilham o mesmo espaço de endereçamento. Isso significa que variáveis globais, memória alocada dinamicamente e outros dados do processo podem ser acessados por todas as threads.

Cada thread, porém, possui elementos próprios, como:

\begin{itemize}
    \item contador de programa;
    \item registradores;
    \item pilha de execução;
    \item estado de escalonamento.
\end{itemize}

Portanto, threads facilitam a comunicação por memória compartilhada, mas exigem cuidado quando múltiplas threads leem e escrevem os mesmos dados.

\subsection{Concorrência}

Concorrência ocorre quando duas ou mais unidades de execução possuem progresso sobreposto no tempo.

Em uma máquina com múltiplos núcleos, esse progresso pode ser efetivamente paralelo. Em uma máquina com um único núcleo, o sistema operacional pode alternar rapidamente entre unidades de execução, produzindo interleaving.

Nos dois programas desta atividade, a concorrência aparece de formas distintas:

\begin{itemize}
    \item no programa com processos, existem quatro processos filhos executando laços de incremento independentes;
    \item no programa com threads, existem quatro threads executando laços de incremento sobre a mesma variável global.
\end{itemize}

\subsection{Condição de corrida}

Uma condição de corrida ocorre quando o resultado de um programa depende da ordem não controlada em que diferentes unidades de execução acessam dados compartilhados.

A operação

\begin{verbatim}
contador++;
\end{verbatim}

parece uma única operação no código-fonte, mas conceitualmente envolve três etapas:

\begin{enumerate}
    \item ler o valor atual de \texttt{contador};
    \item somar uma unidade;
    \item gravar o novo valor em \texttt{contador}.
\end{enumerate}

Se duas threads executarem essas etapas simultaneamente, ambas podem ler o mesmo valor antigo e escrever resultados que sobrescrevem atualizações feitas pela outra thread.

Esse fenômeno é chamado de atualização perdida.

\section{Interpretação didática da tarefa}
\label{sec:atvd03-interpretacao}

A atividade utiliza deliberadamente dois programas pequenos para isolar o conceito observado.

No primeiro programa, o foco é mostrar que a criação de processos não implica compartilhamento automático de variáveis globais. Embora os filhos executem o mesmo código e incrementem uma variável chamada \texttt{contador}, cada um manipula uma cópia privada.

No segundo programa, o foco é mostrar o caso oposto: threads compartilham a variável global, mas esse compartilhamento não garante correção. Como não há sincronização, as atualizações concorrentes sobre \texttt{contador} podem ser perdidas.

A atividade, portanto, não deve ser interpretada como um benchmark de desempenho. O ponto central é semântico:

\begin{center}
\begin{tabular}{|p{0.28\textwidth}|p{0.28\textwidth}|p{0.34\textwidth}|}
\hline
\textbf{Mecanismo} & \textbf{Memória} & \textbf{Efeito sobre \texttt{contador}} \\
\hline
Processos com \texttt{fork()} &
Espaços separados &
Cada filho incrementa sua própria cópia. \\
\hline
Threads POSIX &
Espaço compartilhado &
Todas acessam a mesma variável, mas sem sincronização. \\
\hline
\end{tabular}
\end{center}

Essa diferença é uma base importante para Computação de Alto Desempenho, pois influencia diretamente a forma de decompor problemas, compartilhar dados e preservar a correção.

\section{Modelagem da solução}
\label{sec:atvd03-modelagem}

\subsection{Estrutura geral}

A atividade foi organizada em dois programas independentes.

O primeiro programa, apresentado no Código~\ref{lst:atvd03-cont1}, utiliza
processos. A Figura~\ref{fig:atvd03-modelagem-processos} apresenta essa
modelagem em forma de diagrama UML de componentes.

\begin{figure}[htbp]
    \centering
    \begin{tikzpicture}[
        node distance=0.8cm and 1.0cm,
        every node/.style={font=\small}
    ]
        \node[componente, text width=4.0cm] (pai) {Processo pai\\\texttt{contador = 0}};
        \node[componente, below=of pai, text width=4.0cm] (fork) {\texttt{fork()}\\criação dos filhos};
        \node[processo, below=of fork, text width=5.0cm] (filhos) {4 processos filhos\\cópias privadas de \texttt{contador}\\cada cópia chega a \(1\,000\,000\)};
        \node[componente, below=of filhos, text width=4.0cm] (wait) {\texttt{wait(NULL)}\\sincronização do pai};
        \node[componente, below=of wait, text width=4.0cm] (print) {Imprimir contador\\do processo pai};
        \node[componente, right=1.5cm of filhos, text width=3.4cm] (mem) {Memória isolada\\sem compartilhamento automático};

        \draw[fluxo] (pai) -- (fork);
        \draw[fluxo] (fork) -- (filhos);
        \draw[fluxo] (filhos) -- (wait);
        \draw[fluxo] (wait) -- (print);
        \draw[fluxo] (filhos) -- (mem);
    \end{tikzpicture}
    \caption{Modelagem UML do programa baseado em processos.}
    \label{fig:atvd03-modelagem-processos}
\end{figure}

O segundo programa, apresentado no Código~\ref{lst:atvd03-cont2}, utiliza
threads. A Figura~\ref{fig:atvd03-modelagem-threads} mostra a modelagem
correspondente.

\begin{figure}[htbp]
    \centering
    \begin{tikzpicture}[
        node distance=0.8cm and 1.0cm,
        every node/.style={font=\small}
    ]
        \node[componente, text width=4.0cm] (main) {Processo principal\\\texttt{contador} global};
        \node[componente, below=of main, text width=4.0cm] (create) {\texttt{pthread\_create()}\\criação das threads};
        \node[thread, below=of create, text width=4.7cm] (threads) {4 threads\\executam \texttt{incrementar()}};
        \node[componente, below=of threads, text width=4.7cm] (shared) {Objeto compartilhado\\\texttt{contador++}};
        \node[componente, below=of shared, text width=4.0cm] (join) {\texttt{pthread\_join()}\\impressão final};
        \node[componente, right=1.5cm of shared, text width=3.4cm] (race) {Escrita concorrente\\sem sincronização};

        \draw[fluxo] (main) -- (create);
        \draw[fluxo] (create) -- (threads);
        \draw[fluxo] (threads) -- (shared);
        \draw[fluxo] (shared) -- (join);
        \draw[fluxo] (shared) -- (race);
    \end{tikzpicture}
    \caption{Modelagem UML do programa baseado em threads.}
    \label{fig:atvd03-modelagem-threads}
\end{figure}

Em ambos os casos, as constantes utilizadas são:

\begin{itemize}
    \item quatro unidades de execução;
    \item um milhão de incrementos por unidade.
\end{itemize}

Assim, se todos os incrementos fossem acumulados em um único contador compartilhado e corretamente sincronizado, o valor esperado seria:

\begin{equation}
    4 \times 1\,000\,000
    =
    4\,000\,000.
    \label{eq:atvd03-contador-esperado}
\end{equation}

Esse valor, entretanto, não é o resultado esperado do programa com processos, pois nesse caso o contador não é compartilhado entre pai e filhos.

\subsection{Principais funções e componentes}

Os principais componentes dos programas são apresentados na Tabela~\ref{tab:atvd03-componentes}.

\begin{table}[htbp]
    \centering
    \caption{Principais componentes analisados na Atividade 3.}
    \label{tab:atvd03-componentes}
    \begin{tabular}{p{0.32\textwidth}|p{0.58\textwidth}}
        \hline
        \textbf{Componente} & \textbf{Responsabilidade} \\
        \hline
        \texttt{contador} &
        Variável global usada para evidenciar compartilhamento ou isolamento de memória. \\
        \hline
        \texttt{fork()} &
        Cria processos filhos no primeiro programa. \\
        \hline
        \texttt{wait(NULL)} &
        Faz o processo pai aguardar a finalização dos filhos. \\
        \hline
        \texttt{getpid()} &
        Identifica o processo filho que imprimiu determinado resultado. \\
        \hline
        \texttt{pthread\_create()} &
        Cria threads no segundo programa. \\
        \hline
        \texttt{incrementar()} &
        Rotina executada por cada thread, responsável por incrementar \texttt{contador}. \\
        \hline
        \texttt{pthread\_join()} &
        Faz a thread principal aguardar a finalização das threads criadas. \\
        \hline
    \end{tabular}
\end{table}

\section{Decisões de projeto e trade-offs}
\label{sec:atvd03-decisoes}

\subsection{Uso de variável global}

A variável \texttt{contador} foi declarada globalmente nos dois programas.

Essa escolha torna o experimento mais direto: a mesma forma de declaração produz comportamentos diferentes dependendo do modelo de execução adotado.

No programa com processos, a variável global pertence ao espaço de memória de cada processo.

No programa com threads, a variável global pertence ao processo e é compartilhada pelas threads.

\subsection{Quatro unidades de execução}

O uso de quatro processos e quatro threads facilita a comparação entre os programas.

Em ambos os casos, há quatro unidades concorrentes executando o mesmo trabalho lógico.

Com isso, a diferença observada não está na quantidade de trabalho solicitada, mas na forma como a memória é tratada.

\subsection{Um milhão de incrementos}

O valor de \texttt{INCREMENTOS} foi definido como um milhão.

Esse número é suficientemente grande para tornar perceptível o efeito da condição de corrida no programa com threads. Com poucas iterações, uma execução poderia ocasionalmente produzir um resultado próximo do esperado, mas isso não demonstraria de forma clara o problema.

\subsection{Ausência de sincronização no programa com threads}

O programa com threads não utiliza \texttt{mutex}, \texttt{atomic} ou outra forma de sincronização.

Essa decisão preserva a finalidade didática do código: evidenciar que o compartilhamento de memória exige controle explícito quando há escrita concorrente.

Como trade-off, o resultado final do contador deixa de ser determinístico e não pode ser tratado como correto.

\subsection{Ausência de memória compartilhada explícita no programa com processos}

O programa com processos também não utiliza memória compartilhada, \emph{pipes}, arquivos temporários ou outro mecanismo de comunicação entre processos.

Essa decisão isola o comportamento padrão de \texttt{fork()}: cada processo filho possui sua própria cópia da variável global.

Se o objetivo fosse somar os incrementos dos filhos em um resultado comum, seria necessário acrescentar algum mecanismo de comunicação ou compartilhamento.

\section{Representação estrutural da implementação}
\label{sec:atvd03-estrutura}

A Figura~\ref{fig:atvd03-processos} resume o caso com processos.

\begin{figure}[htbp]
    \centering
    \begin{tikzpicture}[
        node distance=0.8cm and 1.1cm,
        every node/.style={font=\small}
    ]
        \node[componente, text width=3.5cm] (pai) {\textbf{Pai}\\\texttt{contador = 0}};
        \node[processo, right=of pai, text width=4.8cm] (filhos) {\textbf{Filhos 1..4}\\cada processo possui\\\texttt{contador = 1000000}};
        \node[componente, below=of pai, text width=3.5cm] (saida) {\textbf{Saída do pai}\\\texttt{contador = 0}};
        \node[componente, below=of filhos, text width=4.8cm] (isolamento) {Espaços de endereçamento independentes};

        \draw[fluxo] (pai) -- node[above, font=\scriptsize] {\texttt{fork()}} (filhos);
        \draw[fluxo] (pai) -- node[left, font=\scriptsize] {\texttt{wait()}} (saida);
        \draw[fluxo] (filhos) -- node[right, font=\scriptsize] {cópias privadas} (isolamento);
    \end{tikzpicture}
    \caption{Diagrama UML estrutural do programa com processos.}
    \label{fig:atvd03-processos}
\end{figure}

A Figura~\ref{fig:atvd03-threads} resume o caso com threads.

\begin{figure}[htbp]
    \centering
    \begin{tikzpicture}[
        node distance=0.8cm and 1.0cm,
        every node/.style={font=\small}
    ]
        \node[thread, text width=4.4cm] (threads) {\textbf{Threads 1..4}\\ler, somar e gravar};
        \node[componente, right=of threads, text width=4.2cm] (contador) {\textbf{Objeto compartilhado}\\\texttt{contador}};
        \node[componente, below=of contador, text width=4.2cm] (resultado) {Resultado final\\dependente da interleaving};
        \node[componente, below=of threads, text width=4.4cm] (race) {Mesmo processo\\memória compartilhada};

        \draw[fluxo] (threads) -- node[above, font=\scriptsize] {\texttt{contador++}} (contador);
        \draw[fluxo] (contador) -- node[right, font=\scriptsize] {sem sincronização} (resultado);
        \draw[fluxo] (threads) -- (race);
    \end{tikzpicture}
    \caption{Diagrama UML estrutural do programa com threads.}
    \label{fig:atvd03-threads}
\end{figure}

\section{Fluxo de execução}
\label{sec:atvd03-fluxo}

\subsection{Fluxo do programa com processos}

O fluxo do programa \texttt{cont1.c} é:

\begin{figure}[H]
    \centering
    \begin{tikzpicture}[
        node distance=0.55cm and 0.9cm,
        every node/.style={font=\small}
    ]
        \node[processo, text width=2.7cm] (inicio) {Início};
        \node[componente, below=of inicio, text width=4.0cm] (init) {\texttt{contador = 0}\\imprimir valor inicial};
        \node[componente, below=of init, text width=4.0cm] (loop) {Laço de criação\\4 iterações};
        \node[componente, below=of loop, text width=3.2cm] (fork) {\texttt{fork()}};
        \node[decisao, below=of fork, text width=2.4cm] (decide) {É filho?};

        \node[processo, below left=0.9cm and 0.9cm of decide, text width=3.0cm] (filho) {Incrementar contador local};
        \node[processo, below=of filho, text width=3.0cm] (printfilho) {Imprimir PID\\e contador};
        \node[processo, below=of printfilho, text width=3.0cm] (exit) {Encerrar filho};

        \node[componente, below right=0.9cm and 0.9cm of decide, text width=3.0cm] (pai) {Continuar no pai};
        \node[decisao, below=of pai, text width=2.4cm] (more) {4 filhos?};
        \node[componente, below=of more, text width=3.0cm] (wait) {\texttt{wait(NULL)}};
        \node[componente, below=of wait, text width=3.0cm] (printpai) {Imprimir contador\\do pai};
        \node[processo, below=of printpai, text width=2.7cm] (fim) {Fim};

        \draw[fluxo] (inicio) -- (init);
        \draw[fluxo] (init) -- (loop);
        \draw[fluxo] (loop) -- (fork);
        \draw[fluxo] (fork) -- (decide);
        \draw[fluxo] (decide) -- node[above left, font=\scriptsize] {sim} (filho);
        \draw[fluxo] (filho) -- (printfilho);
        \draw[fluxo] (printfilho) -- (exit);
        \draw[fluxo] (decide) -- node[above right, font=\scriptsize] {não} (pai);
        \draw[fluxo] (pai) -- (more);
        \draw[fluxo] (more) -- node[right, font=\scriptsize] {sim} (wait);
        \draw[fluxo] (wait) -- (printpai);
        \draw[fluxo] (printpai) -- (fim);
        \draw[fluxo] (more.east) -- ++(0.7,0) |- node[pos=0.25, right, font=\scriptsize] {não} (fork.east);
    \end{tikzpicture}
    \caption{Diagrama UML de atividade do fluxo com processos.}
    \label{fig:atvd03-fluxo-processos}
\end{figure}

O ponto central desse fluxo é que o pai aguarda os filhos, mas não recebe automaticamente as alterações feitas no contador de cada um.

\subsection{Fluxo do programa com threads}

O fluxo do programa \texttt{cont2.c} é:

\begin{figure}[H]
    \centering
    \begin{tikzpicture}[
        node distance=0.55cm and 0.9cm,
        every node/.style={font=\small}
    ]
        \node[processo, text width=2.7cm] (inicio) {Início};
        \node[componente, below=of inicio, text width=4.2cm] (init) {\texttt{contador = 0}\\imprimir valor inicial};
        \node[componente, below=of init, text width=4.2cm] (create) {Criar 4 threads\\\texttt{pthread\_create()}};
        \node[thread, below=of create, text width=4.2cm] (threads) {Threads 1..4\\executam \texttt{incrementar()}};
        \node[componente, below=of threads, text width=4.2cm] (inc) {Atualizar \texttt{contador++}\\sem sincronização};
        \node[componente, below=of inc, text width=4.2cm] (join) {Aguardar threads\\\texttt{pthread\_join()}};
        \node[componente, below=of join, text width=4.2cm] (print) {Imprimir contador final};
        \node[processo, below=of print, text width=2.7cm] (fim) {Fim};

        \draw[fluxo] (inicio) -- (init);
        \draw[fluxo] (init) -- (create);
        \draw[fluxo] (create) -- (threads);
        \draw[fluxo] (threads) -- (inc);
        \draw[fluxo] (inc) -- node[right, font=\scriptsize] {sem mutex/atomic} (join);
        \draw[fluxo] (join) -- (print);
        \draw[fluxo] (print) -- (fim);
    \end{tikzpicture}
    \caption{Diagrama UML de atividade do fluxo com threads.}
    \label{fig:atvd03-fluxo-threads}
\end{figure}

Nesse caso, o ponto central é que todas as threads escrevem na mesma variável sem sincronização.

\section{Metodologia experimental}
\label{sec:atvd03-metodologia}

\subsection{Ambiente de execução}

Os arquivos da atividade registram os comandos de compilação conceituais:

\begin{verbatim}
gcc cont1.c -o cont1
gcc cont2.c -o cont2 -pthread
\end{verbatim}

O primeiro programa depende das interfaces POSIX \texttt{fork()} e \texttt{wait()}. O segundo depende da biblioteca POSIX Threads.

Nesta consolidação do relatório, foram utilizados os códigos-fonte e as saídas já presentes no diretório da atividade. O objetivo foi documentar o comportamento observado, não refazer um benchmark temporal.

\subsection{Configuração das medições}

Não houve medição de tempo nesta atividade.

A validação experimental consistiu em executar os programas e observar os valores impressos para \texttt{contador}.

Foram considerados:

\begin{itemize}
    \item três execuções registradas do programa com processos;
    \item dez execuções registradas do programa com threads.
\end{itemize}

As execuções do programa com processos foram usadas para confirmar que o contador do processo pai permanece em zero.

As execuções do programa com threads foram usadas para evidenciar que o contador final varia e fica abaixo do valor lógico esperado de \(4\,000\,000\), devido a atualizações perdidas.

\section{Código-fonte desenvolvido}
\label{sec:atvd03-codigo}

\subsection{Programa com processos}

\lstinputlisting[
    style=codigoC,
    caption={Código-fonte da Atividade 3 com processos.},
    label={lst:atvd03-cont1}
]{../ATVD3/cont1.c}

\subsection{Programa com threads}

\lstinputlisting[
    style=codigoC,
    caption={Código-fonte da Atividade 3 com threads.},
    label={lst:atvd03-cont2}
]{../ATVD3/cont2.c}

\section{Validação da implementação}
\label{sec:atvd03-validacao}

\subsection{Validação do programa com processos}

No programa com processos, cada filho realiza exatamente:

\begin{equation}
    1\,000\,000
\end{equation}

incrementos em sua própria cópia de \texttt{contador}.

Assim, espera-se que cada filho imprima:

\begin{equation}
    \texttt{contador} = 1\,000\,000.
\end{equation}

Como o pai não compartilha essa variável com os filhos, espera-se que o processo pai imprima:

\begin{equation}
    \texttt{contador} = 0.
\end{equation}

As três execuções registradas confirmaram esse comportamento.

\subsection{Validação do programa com threads}

No programa com threads, se os incrementos fossem executados de forma atomicamente correta sobre o contador compartilhado, o resultado esperado seria:

\begin{equation}
    4\,000\,000.
\end{equation}

Entretanto, o código utiliza:

\begin{verbatim}
contador++;
\end{verbatim}

sem exclusão mútua ou operação atômica.

Portanto, o resultado final não é deterministicamente igual a \(4\,000\,000\).

As saídas registradas confirmam que o valor final varia entre execuções e permanece abaixo do valor esperado.

\subsection{Observação sobre a impressão inicial}

Nas saídas registradas do programa com processos, a linha

\begin{verbatim}
Contador inicial: 0
\end{verbatim}

aparece mais de uma vez.

Isso não significa que o código tenha executado o trecho inicial várias vezes antes de criar os filhos. O \texttt{printf()} ocorre antes do \texttt{fork()}.

A repetição pode ocorrer quando a saída padrão está sendo redirecionada ou bufferizada. O conteúdo ainda presente no buffer de saída pode ser copiado para os processos filhos durante o \texttt{fork()} e depois descarregado por cada processo.

Esse detalhe reforça outro aspecto importante da criação de processos: além das variáveis, estados associados à biblioteca de E/S também podem ser duplicados no momento da criação do filho.

\section{Resultados experimentais}
\label{sec:atvd03-resultados}

\subsection{Resultados do programa com processos}

Nas três execuções registradas, cada processo filho apresentou contador igual a um milhão, e o processo pai apresentou contador igual a zero.

\begin{table}[htbp]
    \centering
    \caption{Resultados observados no programa com processos.}
    \label{tab:atvd03-processos-resultados}
    \begin{tabular}{c|c|c}
        \hline
        Execução & Contador dos filhos & Contador no pai \\
        \hline
        1 & \(1\,000\,000\) em cada filho & 0 \\
        2 & \(1\,000\,000\) em cada filho & 0 \\
        3 & \(1\,000\,000\) em cada filho & 0 \\
        \hline
    \end{tabular}
\end{table}

O resultado mostra que os filhos executam corretamente seu laço local, mas essas alterações não são refletidas no processo pai.

\subsection{Resultados do programa com threads}

Nas dez execuções registradas, o contador final variou conforme apresentado na Tabela~\ref{tab:atvd03-threads-resultados}.

\begin{table}[htbp]
    \centering
    \caption{Resultados observados no programa com threads.}
    \label{tab:atvd03-threads-resultados}
    \begin{tabular}{c|r}
        \hline
        Execução & Contador final \\
        \hline
        1 & 1\,431\,347 \\
        2 & 1\,305\,740 \\
        3 & 1\,171\,730 \\
        4 & 1\,306\,585 \\
        5 & 1\,235\,595 \\
        6 & 1\,292\,238 \\
        7 & 1\,222\,249 \\
        8 & 1\,681\,956 \\
        9 & 1\,240\,474 \\
        10 & 1\,347\,679 \\
        \hline
    \end{tabular}
\end{table}

O menor valor registrado foi:

\begin{equation}
    1\,171\,730,
\end{equation}

e o maior foi:

\begin{equation}
    1\,681\,956.
\end{equation}

Todos os valores ficaram abaixo de:

\begin{equation}
    4\,000\,000.
\end{equation}

Isso evidencia perda de atualizações durante a execução concorrente.

\section{Visualização dos resultados}
\label{sec:atvd03-visualizacao}

O contraste entre os dois programas pode ser resumido por:

\begin{center}
\fbox{
\begin{minipage}{0.86\textwidth}
\centering
\textbf{Processos}

\vspace{0.2cm}

\[
\text{filho}_i:
\quad
\texttt{contador}=1\,000\,000
\]

\[
\text{pai}:
\quad
\texttt{contador}=0
\]

\vspace{0.2cm}

Memória privada por processo.

\vspace{0.5cm}

\textbf{Threads}

\vspace{0.2cm}

\[
\text{esperado com sincronização}
=
4\,000\,000
\]

\[
\text{observado}
\in
[1\,171\,730,\ 1\,681\,956]
\]

\vspace{0.2cm}

Memória compartilhada com condição de corrida.
\end{minipage}
}
\end{center}

Uma forma textual de visualizar o problema das threads é:

\begin{verbatim}
Valor inicial: contador = 10

Thread A le contador = 10
Thread B le contador = 10

Thread A calcula 11
Thread B calcula 11

Thread A grava 11
Thread B grava 11

Resultado final: 11
Incrementos realizados: 2
Incrementos preservados: 1
\end{verbatim}

Esse exemplo simplificado mostra por que o resultado final pode ficar muito abaixo de quatro milhões.

\section{Análise e discussão}
\label{sec:atvd03-discussao}

\subsection{Comportamento do programa com processos}

O programa \texttt{cont1.c} cria quatro processos filhos.

Cada filho executa um laço com um milhão de incrementos e imprime seu próprio valor de \texttt{contador}.

Como cada filho possui uma cópia independente da variável global, todos os filhos conseguem chegar ao valor:

\begin{equation}
    1\,000\,000.
\end{equation}

O pai, entretanto, nunca incrementa sua própria cópia de \texttt{contador}. Ele apenas cria filhos e aguarda sua finalização.

Portanto, o resultado:

\begin{equation}
    \texttt{contador no processo pai} = 0
\end{equation}

é o comportamento correto para esse código.

Não há, nesse caso, uma condição de corrida sobre \texttt{contador}, porque os processos não estão escrevendo na mesma posição de memória.

\subsection{Problema conceitual no programa com processos}

O problema do primeiro programa não é erro de sincronização.

O problema é de expectativa: se o objetivo fosse obter no pai a soma dos incrementos realizados pelos filhos, o programa atual não fornece comunicação entre processos.

O valor esperado de quatro milhões só faria sentido se os filhos atualizassem um estado comum ou enviassem seus resultados ao pai.

Soluções adequadas seriam:

\begin{itemize}
    \item utilizar memória compartilhada entre processos;
    \item utilizar \emph{pipes} para cada filho enviar seu resultado ao pai;
    \item utilizar arquivos ou outro mecanismo de comunicação;
    \item fazer cada filho retornar uma contribuição e o pai agregá-la por IPC apropriado.
\end{itemize}

Como os filhos são processos independentes, a agregação precisa ser explícita.

\subsection{Comportamento do programa com threads}

O programa \texttt{cont2.c} cria quatro threads dentro do mesmo processo.

Todas executam a função \texttt{incrementar()} e todas acessam a mesma variável global.

Nesse caso, se cada incremento fosse preservado, o resultado deveria ser:

\begin{equation}
    4 \times 1\,000\,000
    =
    4\,000\,000.
\end{equation}

Entretanto, a ausência de sincronização faz com que múltiplas threads executem simultaneamente a sequência leitura-soma-escrita.

Como resultado, parte dos incrementos se perde.

Esse comportamento explica por que as execuções registradas produziram valores como:

\begin{equation}
    1\,431\,347,\quad
    1\,305\,740,\quad
    1\,171\,730.
\end{equation}

Os valores não são iguais entre si porque a ordem de escalonamento das threads pode mudar a cada execução.

\subsection{Problema conceitual no programa com threads}

O problema do segundo programa é uma condição de corrida.

Ao contrário do caso com processos, as threads realmente compartilham a variável. O erro surge justamente porque a escrita concorrente não é controlada.

Soluções adequadas seriam:

\begin{itemize}
    \item proteger o incremento com \texttt{pthread\_mutex\_lock()} e \texttt{pthread\_mutex\_unlock()};
    \item utilizar uma operação atômica para o incremento;
    \item fazer cada thread acumular em uma variável local e somar os quatro resultados ao final;
    \item utilizar primitivas de redução em modelos de programação que as ofereçam, como OpenMP.
\end{itemize}

Para esse caso específico, a solução com acumuladores locais tende a ser mais adequada conceitualmente do que proteger cada incremento individual, pois evita um milhão de sincronizações por thread.

O modelo seria:

\begin{verbatim}
Thread 1 calcula local_1 = 1000000
Thread 2 calcula local_2 = 1000000
Thread 3 calcula local_3 = 1000000
Thread 4 calcula local_4 = 1000000

Thread principal soma:
contador = local_1 + local_2 + local_3 + local_4
\end{verbatim}

Esse padrão preserva a correção e reduz o custo de sincronização.

\subsection{Comparação entre processos e threads}

A Tabela~\ref{tab:atvd03-comparacao} resume as diferenças observadas.

\begin{table}[htbp]
    \centering
    \caption{Comparação conceitual entre processos e threads na atividade.}
    \label{tab:atvd03-comparacao}
    \begin{tabular}{p{0.28\textwidth}|p{0.28\textwidth}|p{0.28\textwidth}}
        \hline
        \textbf{Aspecto} & \textbf{Processos} & \textbf{Threads} \\
        \hline
        Criação &
        \texttt{fork()} &
        \texttt{pthread\_create()} \\
        \hline
        Memória global &
        Copiada para cada processo &
        Compartilhada no mesmo processo \\
        \hline
        Comunicação pelo contador &
        Não ocorre automaticamente &
        Ocorre, mas sem controle \\
        \hline
        Problema principal &
        Falta de comunicação entre processos &
        Condição de corrida \\
        \hline
        Resultado observado &
        Pai permanece com contador 0 &
        Contador final varia e fica abaixo de 4 milhões \\
        \hline
        Solução típica &
        IPC ou memória compartilhada explícita &
        Mutex, operação atômica ou acumulação local \\
        \hline
    \end{tabular}
\end{table}

\section{Limitações}
\label{sec:atvd03-limitacoes}

Os resultados devem ser interpretados dentro do escopo demonstrativo da atividade.

Primeiramente, os programas não medem tempo de execução. Portanto, não é possível concluir, a partir desta atividade, que processos ou threads são mais rápidos para uma aplicação geral.

Em segundo lugar, o programa com threads possui comportamento não determinístico. Os valores registrados demonstram a condição de corrida, mas uma nova execução pode produzir outros valores.

Além disso, o comportamento da impressão inicial no programa com processos pode variar conforme a saída esteja ligada ao terminal ou redirecionada para arquivo. A repetição da linha inicial está associada ao estado do buffer de saída no momento do \texttt{fork()}.

Outra limitação é que o programa com processos não utiliza mecanismos reais de comunicação entre processos. Assim, ele demonstra isolamento de memória, mas não compara diferentes estratégias de IPC.

Por fim, o programa com threads não implementa uma versão corrigida com \texttt{mutex}, operação atômica ou acumuladores privados. A discussão da solução é conceitual, conforme solicitado pelo enunciado.

\section{Conclusões da atividade}
\label{sec:atvd03-conclusao}

A atividade demonstrou que processos e threads oferecem modelos de memória distintos.

No programa com processos, cada filho criado por \texttt{fork()} recebeu uma cópia própria da variável global \texttt{contador}. Por isso, cada filho imprimiu o valor local de um milhão, enquanto o processo pai manteve seu contador igual a zero.

Esse comportamento não representa erro no código. Ele é consequência do isolamento de memória entre processos.

No programa com threads, todas as threads compartilharam a mesma variável global. Nesse caso, a atualização de \texttt{contador} ocorreu sobre um estado comum.

Entretanto, como o incremento não foi protegido por sincronização, ocorreu uma condição de corrida. O resultado final variou entre execuções e ficou abaixo do valor logicamente esperado de quatro milhões.

Assim, a atividade evidencia duas lições complementares:

\begin{enumerate}
    \item processos isolam memória e exigem comunicação explícita quando dados precisam ser combinados;
    \item threads compartilham memória, mas exigem sincronização quando há escrita concorrente.
\end{enumerate}

Em termos de Computação de Alto Desempenho, essa distinção é essencial. A correção de uma aplicação concorrente depende tanto da divisão do trabalho quanto da forma como os dados são compartilhados, protegidos ou combinados.

\section{Síntese}
\label{sec:atvd03-defesa}

Os principais pontos para apresentação oral da atividade são:

\begin{enumerate}
    \item A atividade compara dois programas: um com processos e outro com threads.

    \item Nos dois casos há quatro unidades de execução, e cada uma realiza um milhão de incrementos.

    \item No programa com processos, os filhos são criados com \texttt{fork()}.

    \item Após o \texttt{fork()}, cada processo filho possui sua própria cópia da variável global \texttt{contador}.

    \item Por isso, cada filho imprime \texttt{contador = 1000000}, mas o pai imprime \texttt{contador = 0}.

    \item O pai aguarda os filhos com \texttt{wait(NULL)}, mas aguardar não significa receber as alterações feitas na memória dos filhos.

    \item Se fosse necessário somar os resultados no pai, seria preciso usar comunicação entre processos ou memória compartilhada explícita.

    \item No programa com threads, as threads são criadas com \texttt{pthread\_create()}.

    \item Todas as threads pertencem ao mesmo processo e acessam a mesma variável global.

    \item O valor correto, caso todos os incrementos fossem preservados, seria \(4\,000\,000\).

    \item O código usa \texttt{contador++} sem sincronização.

    \item O incremento não é atomicamente seguro: envolve ler, somar e escrever.

    \item Quando duas threads fazem isso ao mesmo tempo, atualizações podem ser perdidas.

    \item As execuções registradas produziram valores finais variáveis, entre \(1\,171\,730\) e \(1\,681\,956\).

    \item Essa variação caracteriza uma condição de corrida.

    \item Uma solução com threads poderia usar \texttt{mutex}, operação atômica ou acumuladores locais por thread.

    \item Para este problema, acumuladores locais seguidos de soma final são uma solução eficiente, pois evitam sincronizar cada incremento.

    \item A conclusão principal é que processos isolam memória, enquanto threads compartilham memória. Isolamento exige comunicação explícita; compartilhamento exige sincronização.
\end{enumerate}
